\begin{textsample}{POS Dim 1 – am – Score 44.00 – am/am\_ef\_24.txt}  \label{ex:f1_pos_230}
11.
As práticas pedagógicas nas escolas perpassam por uma gama de vivências e estratégias educativas, principalmente no que concerne ao processo de ensino e aprendizagem por meio das \textbf{competências} gerais, que devem ser desenvolvidas de forma integrada aos componentes curriculares, dos objetivos de conhecimento e das habilidades que serão consolidadas pelos estudantes ao longo das etapas/modalidades da Educação Básica, evidenciadas neste Referencial Curricular Amazonense aliado à Base Nacional Comum Curricular = BNCC.
Destaca -se que todo processo educativo \textbf{demanda} como ponto estratégico ferramentas essenciais para a verificação do aprendizado, as avaliações internas e externas à escola.
A avaliação interna, mais conhecida como avaliação da aprendizagem, de natureza processual e formativa, elaborada e realizada pelo professor da sala de \textbf{aula}, busca \textbf{verificar} a aprendizagem dos estudantes e é produzida em conformidade com o planejamento escolar.
Esse tipo de avaliação requer que o professor diversifique os diferentes instrumentos \textbf{avaliativos}, com base nas habilidades abarcadas no processo de construção de determinados conhecimentos, com a finalidade de \textbf{subsidiar} \textbf{tanto} o trabalho pedagógico quanto o acompanhamento individual no desenvolvimento dos estudantes em suas aprendizagens essenciais \textbf{esperadas} para cada ano, uma vez que integra as habilidades do currículo.
No que tange à avaliação externa, em larga escala, elaborada e realizada por agentes externos à escola, aplicada de forma abrangente a uma rede de ensino municipal, estadual, privada, ou em várias redes de ensino, é construída por meio de testes padronizados, com base nas \textbf{competências} definidas nas \textbf{Matrizes} de Referências, que é um recorte do currículo, com a finalidade de identificar o nível de desempenho de escolas e das redes de ensino, \textbf{verificando} se os estudantes aprenderam o que de fato é proeminente para cada etapa/ano de ensino avaliado.
É uma ferramenta que fornece elementos para a formulação e o monitoramento de políticas públicas, bem como o redirecionamento de práticas pedagógicas e de tomada de decisões.
Sendo assim, embora as duas avaliações tenham papéis diferenciados, elas se complementam, pois, para que o estudante consiga desenvolver tais \textbf{competências} cobradas nas avaliações externas de larga escala, o professor precisará desenvolver cotidianamente \textbf{inúmeras} habilidades \textbf{agregadas} a ela.
Dessa forma, a realização de cada \textbf{aula}, proporcionará melhor aprendizagem, portanto melhor desempenho das \textbf{competências} e habilidades por parte dos estudantes em cada componente curricular.
Contudo, é importante que os profissionais de educação no âmbito escolar estabeleçam as relações entre os indicadores obtidos nas avaliações externas e o rendimento apresentado pelos estudantes no processo de \textbf{ensino-aprendizagem}, isto é, nas avaliações internas, realizada pelos professores no dia a dia da sala de \textbf{aula}.
Diante do \textbf{exposto}, a avaliação é, portanto, parte integrante do currículo, na medida em que a ele se congrega como uma das etapas do processo pedagógico.
Assim, a avaliação, em sentido amplo, não deve ser classificatória e excludente, mas uma prática pedagógica que possibilite a \textbf{análise} e reflexão do processo de ensino e aprendizagem, possibilitando novas estratégias de ensinar e aprender, tornando -se ferramenta \textbf{imprescindível} das políticas públicas educacionais e da potencialização da aprendizagem que todo estudante deve desenvolver, articulado aos variados conhecimentos que circundam o espaço e cotidiano escolar.
As Secretarias de Educação devem contribuir a partir da efetivação de políticas educacionais igualitárias, efetivas e impactantes, capazes de melhorar significativamente a aprendizagem dos estudantes, pois, o Referencial Curricular Amazonense e a BNCC em consonância com os princípios democráticos e inclusivos visam à formação global dos estudantes para que se tornem de fato competentes diante dos desafios apresentados pela sociedade \textbf{contemporânea}.
Com isso, reforça -se a prática \textbf{avaliativa} como um processo natural e fundamental para a qualidade da educação brasileira.
Importa \textbf{dizer} que a prática da avaliação dos sistemas de ensino ocorreu gradualmente, mas somente \textbf{atingiu} seu marco legal a partir da reforma educacional na década de 90, em que o governo passou a incorporar metodologias de gestão voltadas para a melhoria dos resultados dos serviços prestados à população.
De acordo com Depresbiteris ( 2001, p.144 ), existem três propósitos da avaliação nos sistemas de ensino : > “ Fornecer resultados para a gestão da educação ; \textbf{subsidiar} a melhoria dos projetos pedagógicos das escolas e propiciar informações para a melhoria da própria avaliação, o que a caracteriza como meta-avaliação ”.
É por meio desse instrumento que são produzidos indicadores comparativos de desempenho, que servirão de base para o monitoramento, a ( re ) formulação de políticas públicas, assim como da gestão da educação, no âmbito da escola e nas diferentes esferas do sistema educacional.
A LDBEN Nº 9.394/96, Art. 9º aponta a necessidade de um controle por parte do Estado Brasileiro desde o nível básico de ensino até o superior, por meio das avaliações centralizadas na União, que são \textbf{explicitadas} no Plano Nacional de Educação — PNE, em vigência, como parte do Sistema Nacional de Avaliação da Educação Básica - SAEB.
Com isso, a Educação passou a ser ajustada como um serviço que precisava ser gerido de acordo com os resultados apresentados, ou melhor, em consonância com a"gestão de resultados ”.
Para inserir o país na vida \textbf{moderna} \textbf{globalizada}, assim como os demais países, o Estado necessitava criar uma política de avaliação nacional de acompanhamento do trabalho administrativo e pedagógico das escolas para responder às expectativas internacionais.
A partir da CF/1988, a avaliação da educação básica passa a ser necessária para a obtenção da qualidade tão difundida pela política neoliberal, pois em seu Art. 206 dispõe sobre a “ garantia de qualidade ”.
No Art. 209, a referida Constituição versa sobre a “ autorização e avaliação de qualidade pelo poder público ”, e em seu Art. 214 que trata do PNE fala a respeito da “ melhoria da qualidade do ensino ”.
Assim, por meio do Ministério da Educação - MEC foi criado um abrangente e \textbf{complexo} sistema de avaliação no Brasil.
Em meados dos anos 90, impulsionado pela Lei maior do Estado, compromissos internacionais e pelos processos \textbf{avaliativos} realizados anteriormente nos estados e municípios, o MEC em articulação com as Secretarias Estaduais de Educação implantou o SAEB, levando as avaliações externas em larga escala do país se destacarem em âmbito nacional.
Desde então, o MEC tem buscado realizar, periodicamente, a aplicação da avaliação do desempenho escolar dos estudantes, com o \textbf{intuito} de melhorar a qualidade do ensino nas escolas, já que seu principal objetivo, de acordo com o Instituto Nacional de Estudos e Pesquisas Anísio Teixeira - INEP, é avaliar a Educação Básica brasileira, contribuir para a melhoria de sua qualidade e para a universalização do acesso à escola, oferecendo \textbf{subsídios} concretos para a formulação, reformulação e o monitoramento das políticas públicas voltadas para a Educação Básica.
Em \textbf{vista} disso, é notória a implantação de sistemas próprios de avaliação externa em larga escala pelas redes de ensino em nível estadual e municipal no \textbf{decorrer} das últimas décadas, com o \textbf{intuito} de melhorar a \textbf{proficiência} dos estudantes e a qualidade da educação básica em âmbito local.
Cumprir essa \textbf{tarefa} envolve desafios, como enfrentar as desigualdades extras e intraescolares : a pobreza e a violência ; as novas formas de estrutura familiar ; as particularidades de cada local, de cada escola e do desenvolvimento cognitivo de cada estudante, entre outras questões.
Diante disso, é necessário reunir informações concretas sobre a população atendida e o ensino ofertado para, desse modo, implementar ações que visem \textbf{atingir} o objetivo traçado.
Nessa perspectiva, a educação básica ganhou maior visibilidade frente à sociedade que começou a avaliar o nível de ensino por meio das avaliações externas aplicadas, que além de \textbf{verificar} se o currículo oficial está sendo cumprido, possibilita também a elaboração de políticas públicas aos sistemas de ensino.
Destarte, a utilização da avaliação educacional no Brasil com \textbf{foco} nos resultados faz parte das diretrizes do Plano de Metas Compromisso Todos pela Educação, implementado em 2007 sob a orientação do Plano de Desenvolvimento da Educação - PDE para toda a rede de ensino, caracterizado conforme o Art. 1º do Decreto Nº 6.094, de 24 de abril de 2007.
Tal plano foi uma estratégia do governo de \textbf{mobilizar} a sociedade para a efetivação das metas através da adesão dos Estados e municípios.
Com isso, haveria maior possibilidade de oferecer um atendimento de qualidade aos educandos da rede pública de ensino com o intenso trabalho das escolas em alcançar as metas propostas pelo Índice de Desenvolvimento da Educação Básica - IDEB, que permite"identificar quais são as redes de ensino estaduais, municipais e as escolas que apresentam maiores \textbf{fragilidades} no desempenho escolar e que, por isso mesmo necessitam de maior atenção e apoio financeiro e de gestão ” ( BRASIL, 2011, p. 4 ).
O IDEB é um dos indicadores nacionais construídos a partir dos resultados das avaliações externas em larga escala como a Avaliação Nacional do Rendimento Escolar ANRESC, mais conhecida como Prova Brasil. > Os resultados do Saeb e da Prova Brasil são importantes, pois contribuem para dimensionar os \textbf{problemas} da educação básica brasileira e orientar a formulação, a implementação e a avaliação de políticas públicas educacionais que conduzam à formação de uma escola de qualidade ( BRASIL, 2011, p.5 ).
Essas avaliações por meio de seus resultados fornecem \textbf{subsídios} para a tomada de decisões destinadas a benfeitorias no sistema de ensino e nas escolas através do acompanhamento de suas diferentes \textbf{edições}.
De acordo com Castro ( 2009 ), a avaliação externa em larga escala é necessária : > Independente dos motivos que levam à criação de sistemas de avaliação, parece haver concordância quanto ao seu importante papel como instrumento de melhoria da qualidade.
Como os resultados da educação não são diretamente observáveis nem \textbf{imediatos}, dada a heterogeneidade do corpo docente e da situação socioeconômica familiar dos \textbf{alunos}, só é possível obter uma visão geral do desempenho dos sistemas educacionais mediante uma avaliação externa em larga escala ( CASTRO, 2009, p. 6 ).
Dessa forma, a política de avaliação educacional, estruturada no SAEB até o ano de 2017, fez uso de avaliações de desempenho dos estudantes para aferir a qualidade da educação brasileira nas etapas do Ensino Fundamental e \textbf{médio} por meio de avaliações como : Provinha Brasil, Avaliação Nacional da Alfabetização ( ANA ), ANRESC/Prova Brasil ; o \textbf{Exame} Nacional do Ensino Médio ( ENEM ), o \textbf{Exame} Nacional de Certificação de Jovens e Adultos ( ENCCEJA ) e o Programa Internacional de Avaliação de Estudantes ( PISA ).
As referidas avaliações nacionais permitiram ao MEC e às Secretarias Estaduais e Municipais de Educação, a definição de ações para solucionar os \textbf{problemas} identificados, direcionar os recursos técnicos e financeiros aos setores mais necessitados e \textbf{fomentar} o \textbf{debate} acerca do trabalho pedagógico exercido nas diversas instituições escolares sob a sua jurisdição, a fim de desenvolver mais eficazmente o sistema educacional brasileiro.
A partir da BNCC, o SAEB vem sendo \textbf{aprimorado}, e novos \textbf{encaminhamentos} por parte do MEC/INEP são socializados com as redes de ensino quanto ao abandono das siglas e do nome fantasia das avaliações, para a unificação da nomenclatura SAEB para todas as avaliações externas em larga escala a partir de 2019.
O sistema contemplará todas as etapas de ensino, pois passará a incluir a Educação Infantil, juntamente com o Ensino Fundamental e Médio.
As aplicações acontecerão para turmas de creche e pré-escola, 2º ano, 5º ano, 9º ano e 3º série do Ensino Médio.
Vale ressaltar que o 3º ano do Ensino Fundamental não será mais avaliado, tendo em \textbf{vista} as diretrizes da BNCC que estabelece a consolidação do ciclo da alfabetização ao \textbf{final} do 2º ano.
Quanto ao 9º ano, além dos componentes curriculares de Língua Portuguesa e Matemática, serão contempladas também no teste padronizado as áreas de \textbf{Ciências} da Natureza e Ciências Humanas.

% matched lemmas: agregar, aluno, análise, aprimorar, atingir, aula, avaliativo, ciência, competência, complexo, contemporâneo, debate, decorrer, demanda, dizer, edição, encaminhamento, ensino-aprendizagem, esperar, exame, explicitar, exposto, final, foco, fomentar, fragilidade, globalizado, imediato, imprescindível, intuito, inúmero, matriz, mobilizar, moderno, médio, problema, proficiência, subsidiar, subsídio, tanto, tarefa, verificar, vista
\end{textsample}
